\section{Introduction}
\label{sec:intro}

A model that scores 83\% on a bias benchmark can score 99\% on the same questions stripped of their answer options.  The model has not changed; the format has.  This paper is about what changes when the evaluation format changes, and about what happens to that format when an agentic scaffold restructures the query before the model ever sees it.  The two questions turn out to be the same question.

\paragraph{Two gaps.}
The first is architectural.  Almost every major benchmark cited in safety evaluations~\cite{lin2022truthfulqa, parrish2022bbq, rottger2024xstest} and almost every responsible-scaling framework~\cite{anthropicrsp2026} and almost every third-party assessment examines a model as a standalone function: prompt in, response out, score.  Real-world production models are not standalone functions.  They sit inside agentic scaffolds that add reasoning traces~\cite{yao2023react}, run them through critic agents~\cite{du2024improving}, or decompose tasks across delegation pipelines.  Each restructures the input, and therefore changes, in principle, what the safety benchmark is measuring.
The second gap is representational.  Most widely-used safety benchmarks are multiple-choice~\cite{lin2022truthfulqa, parrish2022bbq, perez2023discovering}; scaffolded and agentic deployments interact with models through open-ended natural language.  A growing literature has documented format sensitivity for capability evaluation: position bias in MC selection~\cite{zheng2024selectionbias}, 13--75\% accuracy swings from answer reordering~\cite{pezeshkpour2024ordersensitivity}, ${\sim}25$~pp drops moving from MC to OE~\cite{myrzakhan2024openllm}.  Neither gap is news on its own: format sensitivity is established for capability evaluation, and task decomposition obviously strips context.  The effects are real but oddly shaped: format effects on safety run in property-dependent directions, with sign reversal across the four constructs studied; the within-format scaffold contrast collapses to a near-null; and 35.6~pp of opposing-sign model$\times$scaffold movement happens on the same items.

\paragraph{What this paper contributes.}
The paper does three things.  It quantifies what the two gaps actually contribute to safety assessment (5--20~pp shifts on identical items, more than enough to defeat most single-factor explanations).  It collapses the scaffold finding and the format finding into the same mechanism: map-reduce significantly reduces measured safety by modifying the evaluation format mid-flight (decomposition eliminates MC options from sub-calls; the resulting effective OE state is scored against an MC key), and once that channel is closed the residual is small, model-varied, and reflects real reasoning disruption from delegation rather than format conversion.  And it lays out the interaction structure (model$\times$scaffold spans of 35.6~pp, pooled flatness inside a fixed format, per-property directionality) that any policy framework conditioning deployment on safety scores has to accommodate: only per-configuration testing yields useful safety signal.

\paragraph{What we did.}
The paper has four pre-registered components (three pre-registered, one pre-data frozen); items, scoring, and registration status vary across them.  The consolidated dataset mapping (Table~\ref{tab:dataset-map}) at the start of Section~\ref{sec:overview} identifies which sample backs which estimate.  The main scaffold study ($N = 62{,}808$; six frontier models, four deployment configurations) was pre-registered on OSF (DOI:~10.17605/OSF.IO/CJW92).  It evaluates four safety benchmarks (BBQ, TruthfulQA, XSTest/OR-Bench, sycophancy; the sycophancy benchmark was added by addendum after the originally registered items proved to be a non-safety property, retained as the AI factual-recall control study below), with blinded scoring, equivalence testing~\cite{schuirmann1987tost, lakens2017equivalence}, and a specification curve enumerating 29 researcher degrees of freedom (the primary curve passes through 3 of these via 18 specifications; a broader exploratory curve passes through 9 via 384).  A separately pre-registered Phase~2 mechanistic component ($N = 7{,}200$; DOI:~10.17605/OSF.IO/WA9Y7) examines four levels of invocation intensity on a sample disjoint from the primary set.  The AI factual-recall control ($N = 12{,}000$) repurposes the items the original CJW92 pre-registration mis-labelled as sycophancy: a non-safety within-paper null on which scaffold effects should be near zero if scaffolds disrupt task-following rather than safety per se.  The format-dependence experiment ($N = 4{,}400$; five models, five metrics) was added after the primary scaffold results implicated format conversion rather than misalignment as the driver of the map-reduce effect; its design was frozen before data collection, but its status is pre-data exploratory rather than ab initio confirmatory, and the paper labels it accordingly throughout.

\paragraph{What we found.}
Two of three scaffold architectures preserve safety within practically meaningful pooled margins (ReAct RD$\,=\,-0.7$~pp, statistically significant in the full sample but TOST-equivalent within $\pm 2$~pp and not significant after excluding Opus; multi-agent RD$\,=\,-0.6$~pp, non-significant, TOST-equivalent).  Map-reduce delegation produces the substantial pooled degradation (RD$\,=\,-7.3$~pp, NNH$\,=\,14$).  Chasing that result revealed a deeper problem: on identical items, MC vs.\ OE administration shifts safety scores by 5--20~pp.  Within-format scaffold contrasts are pooled-flat (cell-level CIs descriptive at our cell sizes); format conversion, not scaffold architecture, emerges as the operative variable for the bulk of the map-reduce loss.  An option-preserving variant recovers 40--89\% of that loss on the two benchmarks where we ran it; the residual 11--60\% is real and varies across models.

Of the four safety properties tested, sycophancy has by far the lowest baseline (29.2\% non-sycophantic at direct API, 31.0\% pooled across the four configurations).  All three scaffolds improve it in aggregate, and the model-by-scaffold dispersion is the widest in the study: on the same items, under map-reduce, Opus loses 16.8~pp while Llama~4 gains 18.8~pp.

Variance decomposition arrives at the same picture from another angle: scaffold architecture accounts for $0.4\%$ of outcome variance, while benchmark choice accounts for roughly $45\times$ more (19.3\%).  The generalizability analysis returns $G = 0.000$ with bootstrap 95\% CI $[0.000, 0.752]$.  Composite reliability cannot be distinguished from zero on the four-benchmark mix, and the upper bound does not rule out moderate reliability under a richer one.  The width of the interval alone defeats a single composite safety number as a deployment input.

\paragraph{Implications.}
Current responsible-scaling policies condition deployment decisions on direct-API benchmark scores.  Those scores are format-contingent measurements that shift substantially with administration format, and structure-destroying scaffolds change the effective administration format without the evaluator's awareness.  We restrict scope deliberately to proxy safety properties (bias, sycophancy, truthfulness, over-refusal): if even these well-understood proxies exhibit format sensitivity and scaffold fragility, the field's ability to reliably evaluate catastrophic risks under agentic deployment is, in the most generous reading, an open question.  Three concrete mandates for evaluation reform follow in Section~\ref{sec:policy}.

\paragraph{A structural limitation up front.}
The evaluation pipeline was built primarily using Claude Opus~4.6 (with GPT and Gemini at various development stages).  Opus~4.6 is one of the six models being tested.  Section~\ref{sec:reflexivity} enumerates the threat channels arising from shared use of Opus in both direct-API testing and pipeline-builder functions, along with structural safeguards (pre-registration, fully automated scoring, blinded assessor, shared LiteLLM code path without model-specific branching), the Opus-excluded sensitivity analysis (which maintains the findings related to map-reduce but demotes the small ReAct effect to non-significance), and the remaining risks (subtler bias channels that aggregate hypothesis tests cannot detect).  The honest framing is that Opus's combined direct-API and pipeline-builder advantage warrants independent replication with a pipeline developed on a different model; the source code is released so that such replication is mechanically possible.
